\section{Ratio analysis}

These entries describe the tools of financial analysis and the main ratio families, use relationships among ratios to evaluate a company (including industry-specific measures), decompose return on equity and attribute its change, and show how ratios drive an earnings forecast.

\begin{qfmethod}[Tools of financial analysis and the activity, liquidity, solvency and profitability ratio families: definitions, calculation and interpretation]\label{meth:financial-ratio-families}
\qffield{Statement} \emph{Tools} \citep[pp.~28--35]{financial-statement-analysis-2009}:
\begin{itemize}
\item comparative (horizontal) analysis of year-to-year changes and index trends;
\item common-size (vertical) statements;
\item ratio analysis, a starting point that identifies questions;
\item cash-flow analysis;
\item valuation.
\end{itemize}
Ratios need an economically meaningful relation between numerator and denominator. They are interpreted against history, standards and competitors, and are affected by economic events, industry, management policy and accounting methods (pp.~34--35).

\emph{Families} (\hyperref[ex:harbour-tools-statements]{Harbour}, using average balances for flow-to-stock ratios):

\begin{center}\small
\begin{tabular}{@{}llrr@{}}
\toprule
Family & Ratio & 2025 & 2024\\
\midrule
Activity & Days sales outstanding (revenue$/$avg.\ receivables) & 58.3 & 57.5\\
 & Days of inventory (cost of sales$/$avg.\ inventory) & 107.9 & 108.4\\
 & Days payables (purchases$/$avg.\ payables) & 91.1 & 92.5\\
 & Cash conversion cycle (days) & 75.1 & 73.4\\
 & Fixed-asset turnover (revenue$/$avg.\ net PP\&E) & 1.85 & 1.81\\
 & Total-asset turnover & 0.874 & 0.846\\
Liquidity & Current ratio & 1.84 & 1.87\\
 & Quick ratio & 0.93 & 0.94\\
Solvency & Debt to equity & 0.55 & 0.55\\
 & Net debt to EBITDA & 1.63 & 1.50\\
 & EBIT interest cover & 5.76 & 6.11\\
Profitability & Gross margin (\%) & 35.46 & 36.00\\
 & Operating margin (\%) & 8.61 & 11.00\\
 & Return on assets (\%) & 4.48 & 5.66\\
 & Return on equity (\%) & 8.84 & 11.52\\
\bottomrule
\end{tabular}
\end{center}

\qffield{Notation} Days $=365/$turnover. Cash conversion cycle $=$ days sales outstanding $+$ days of inventory $-$ days payables. Payables turnover uses purchases on credit from suppliers (Harbour: 485 in 2025). Cost of sales plus the change in inventory overstates them when inventory also absorbs labour and depreciation; the book's version uses cost of sales (p.~462).

\qffield{Conditions} Averages smooth year-end effects but not seasonality; companies with a natural business year (for example retailers at a low-inventory year-end) need interim averages (p.~449). Ratios computed on unadjusted statements inherit accounting differences (\hyperref[meth:accounting-analysis-adjustments]{adjustments}).

\qffield{Practice} \emph{Interpretation} (Harbour).
\begin{itemize}
\item Activity was broadly stable, with slightly slower collection and faster payment lengthening the cycle by 1.7 days.
\item Asset turnover improved as revenue grew faster than assets.
\item Liquidity is adequate but inventory-dependent.
\item Leverage measured on EBITDA rose because EBITDA fell, and on equity it was flat.
\item Profitability fell on every measure, mainly through the operating margin.
\end{itemize}
The next entry evaluates these relationships.

\qffield{Pitfalls} Computing many overlapping ratios; ratios share components, so a few well-chosen ones suffice (p.~35). Mixing year-end and average denominators across years. Using cost of sales for payables turnover in a company with large inventory changes.

\qffield{See also} \hyperref[meth:ratio-based-evaluation]{evaluating with ratios}; \hyperref[meth:balance-sheet-ratios]{balance-sheet ratios}; \hyperref[meth:income-statement-ratios]{income-statement ratios}.

\end{qfmethod}

\begin{qfmethod}[Relationships among ratios and evaluating a company with them, including industry-specific ratios and their uses and limits]\label{meth:ratio-based-evaluation}
\qffield{Statement} Ratios are linked, so evaluation follows their relationships rather than a list. The main links are:
\begin{itemize}
\item return on assets is margin times turnover;
\item return on equity adds leverage (\hyperref[meth:return-on-equity-decomposition]{decomposition});
\item the cash conversion cycle links activity to liquidity and financing needs;
\item coverage links profitability to solvency;
\item liquidity ratios depend on inventory and receivable quality.
\end{itemize}
Evaluation compares each against the company's history, its competitors and the analyst's criteria, with attention to variability as well as trend \citep[p.~35]{financial-statement-analysis-2009}. A company's building blocks are interrelated: operating performance affects liquidity and financing, and credit standing depends on profitability as well as liquidity (p.~638).

\emph{Industry-specific measures.} Industries with special accounting or regulation need specialist knowledge and measures. Examples are oil and gas cost centres and depletion, insurance premium revenue and reserves, and utilities' rate base (p.~640). Sector-specific valuation measures divide enterprise value by operating units such as tons of capacity, kilowatt hours, barrels of reserves or subscribers \citep[pp.~571--576]{investment-valuation-2012}. They link value to operations, can be computed where accounts are unreliable or not comparable, and are sometimes used ``in desperation''. Their limits are that they can let a whole sector look cheap relative to itself and make fundamental differences hard to control for. Comparisons need adjustment for each unit's cash generation, acquisition cost and growth.

\qffield{Notation} ROA $=$ margin $\times$ turnover. Value per unit $=$ enterprise value$/$number of units.

\qffield{Conditions} Industry-specific measures are meaningful only within an industry with homogeneous units; they are not substitutes for cash-flow analysis.

\qffield{Practice} \emph{Evaluation} (\hyperref[ex:harbour-tools-statements]{Harbour}, 2025 against 2024).
\begin{itemize}
\item \emph{Profitability.} Return on assets fell from 5.66 to 4.48~percent. Turnover rose (0.846 to 0.874), so the fall is a margin problem: gross margin $-0.54$ points from cost inflation not passed through, plus restructuring.
\item \emph{Activity and liquidity.} The cash conversion cycle is stable at about 75 days. The quick ratio below one is normal for a manufacturer with 108 days of inventory, but makes liquidity depend on selling stock.
\item \emph{Solvency.} Net debt to EBITDA rose to 1.63 (1.49 excluding restructuring) and EBIT cover fell to 5.8 times. These remain moderate, but the year-end RCF drawing shows investment and dividends exceed internal funds.
\end{itemize}
\emph{Conclusion.} Harbour is sound but under margin pressure. The decisive questions are pricing power and the payoff of the restructuring, which the forecast tests.

An industry measure for a tool maker, such as enterprise value per unit of annual production capacity, would add little here. Units are heterogeneous, and the accounts are reliable.

\qffield{Pitfalls} Evaluating ratios one at a time. Treating sector-specific multiples as fundamental value. Ignoring that improved turnover can come from under-investment.

\qffield{See also} \hyperref[meth:financial-ratio-families]{ratio families}; \hyperref[def:statement-analysis-framework]{framework}.

\end{qfmethod}

\begin{qfmethod}[DuPont analysis of return on equity: three- and five-factor decompositions, operating and financing components, and attributing a change to its drivers]\label{meth:return-on-equity-decomposition}
\qffield{Statement} Return on common equity (ROE) is profit to ordinary shareholders over average common equity \citep[p.~449]{financial-statement-analysis-2009}. It decomposes as follows.
\begin{itemize}
\item Three factors: net margin $\times$ asset turnover $\times$ leverage (average assets$/$average equity) (p.~463).
\item Five factors, splitting net margin: tax burden (profit$/$PBT) $\times$ interest burden (PBT$/$EBIT) $\times$ EBIT margin $\times$ asset turnover $\times$ leverage.
\item Operating and financing parts: $\text{ROE}=\text{RNOA}+\text{FLEV}\times(\text{RNOA}-\text{NFR})$ (pp.~463--465). Return on net operating assets (RNOA) is NOPAT over average net operating assets. FLEV is average net financial obligations over average equity. The net financial rate (NFR) is after-tax net financial expense over average net financial obligations.
\end{itemize}
Leverage raises ROE only while the spread is positive. Operating returns are more persistent and more highly valued than returns from leverage (p.~464).

\qffield{Notation} The identities are exact when the same averages are used throughout. Exact attribution of a change uses logs: $\ln(\text{ROE}_1/\text{ROE}_0)=\sum_k\ln(f_{k,1}/f_{k,0})$.

\qffield{Conditions} Here operating tax is total tax plus the statutory shield on finance costs, so that NOPAT less after-tax net financial expense equals profit exactly. Harbour's net financial obligations include lease liabilities and the pension deficit, consistent with their interest in finance costs; including or excluding them changes the split, not ROE.

\qffield{Practice} \emph{Original} (\hyperref[ex:harbour-tools-statements]{Harbour}).

\begin{center}\small
\begin{tabular}{@{}lrrr@{}}
\toprule
 & 2025 & 2024 & Contribution to change (points)\\
\midrule
Tax burden & 0.7207 & 0.7283 & $-0.105$\\
Interest burden & 0.8262 & 0.8364 & $-0.123$\\
EBIT margin & 0.0861 & 0.1100 & $-2.479$\\
Asset turnover & 0.8741 & 0.8455 & $+0.337$\\
Leverage & 1.9731 & 2.0340 & $-0.308$\\
ROE & 8.84\% & 11.52\% & $-2.678$\\
\bottomrule
\end{tabular}
\end{center}
The log contributions sum exactly to the change. The EBIT margin explains 2.48 of the 2.68-point fall. Restructuring alone is 1.67 points of margin; excluding it, ROE would be 11.00~percent.

Operating and financing split: RNOA was 6.85~percent (8.53 in 2024), the net financial rate 3.37~percent and FLEV 0.574. The spread of 3.48 points adds $0.574\times3.48=2.00$ points, so ROE is 8.84~percent. About 77~percent of ROE comes from operations, against 74~percent in 2024. The book reports an average of about 84~percent for public companies \citep[p.~465]{financial-statement-analysis-2009}, citing Nissim and Penman (2001), which was not read here.

\qffield{Pitfalls} Reading a higher ROE from more leverage as better performance. Mixing year-end and average denominators so that factors do not multiply to ROE. Attributing changes with sequential substitution in an arbitrary order and treating the result as unique.

\qffield{See also} \hyperref[meth:financial-ratio-families]{ratio families}; \hyperref[meth:ratio-based-earnings-forecast]{forecasting}.

\end{qfmethod}

\begin{qfmethod}[Using ratio analysis to model and forecast earnings: driver ratios, earning power, disaggregation, reasonableness checks and sensitivity]\label{meth:ratio-based-earnings-forecast}
\qffield{Statement} An earnings forecast projects the components, not the total. The book's procedure \citep[pp.~492--498]{financial-statement-analysis-2009}:
\begin{itemize}
\item start from sales growth informed by the economy, competition and expansion plans;
\item apply margin and expense ratios, either recent or adjusted for expected conditions;
\item project depreciation from the asset base, interest from debt and tax from the rate;
\item derive balance-sheet items from turnover ratios.
\end{itemize}
Forecasting differs from extrapolation. Earnings are analysed by segment when segments differ in risk, profitability or growth, and the forecast is checked for reasonableness against past and industry returns on invested capital (pp.~619--620). Earning power, the level expected to persist, is best measured over several years (p.~617). Sensitivity analysis varies the assumptions with the largest effects (pp.~498--499).

\qffield{Notation} Profit $=$ revenue $\times$ EBIT margin $-$ interest $-$ tax. Sensitivity of profit to the EBIT margin $\approx$ revenue $\times(1-t)$, before second-round interest effects.

\qffield{Conditions} Ratio-driven forecasts assume stable relations; when the business or pricing environment changes (\hyperref[meth:competition-inflation-forecasting]{competition and inflation}), the ratios must be re-derived from drivers.

\qffield{Practice} \emph{Original} (\hyperref[ex:harbour-pro-forma-model]{Harbour forecast}, 2026).
\begin{itemize}
\item Revenue grows 6.09~percent (3~percent volume and 3~percent price) to 1{,}145.77.
\item Gross margin is 35.29~percent, with 4~percent variable-cost and 3~percent fixed-cost inflation.
\item Selling and administration is 17.5~percent of revenue in cash, plus the share-based, pension and right-of-use charges.
\item EBIT is 119.52 (a 10.43~percent margin, against 8.61 in 2025 and 11.00 in 2024), finance costs 18.50 and profit 73.40.
\end{itemize}
\emph{Reasonableness.} The implied RNOA of 8.39~percent sits between 2025 (6.85) and 2024 (8.53). The forecast is a partial margin recovery once restructuring ends, not a step change. \emph{Sensitivity.} A one-point lower EBIT margin cuts 2026 profit by 8.76 in the full model, against $1{,}145.77\times1\%\times0.75=8.59$ by the approximation; the gap is higher revolver interest.

\qffield{Pitfalls} Extrapolating last year's margin through a known transitory item. Holding asset turnover fixed when growth requires capacity. Skipping the return-on-capital check.

\qffield{See also} \hyperref[ex:harbour-pro-forma-model]{pro forma model}; \hyperref[meth:earnings-quality-persistence]{persistence}; \hyperref[meth:forecast-biases-remedies]{forecast biases}.

\end{qfmethod}
